\documentclass[8pt, a4paper, landscape]{extarticle}
% Компактные поля
\usepackage[margin=1cm]{geometry}
\usepackage[T2A]{fontenc}
\usepackage[utf8]{inputenc}
\usepackage[english, russian]{babel}
\usepackage{amsmath, amssymb, amsfonts}
\usepackage{enumitem}
\usepackage{multicol} % Для колонок
\usepackage[most]{tcolorbox}

% Настройка списка для экономии места
\setlist[itemize]{leftmargin=*, noitemsep, topsep=0pt}
\setlist[enumerate]{leftmargin=*, noitemsep, topsep=0pt}

% Стиль блока билета
\newtcolorbox{ticket}[1]{
  colback=white, colframe=black!70, coltitle=white,
  fonttitle=\bfseries\small, title={#1}, arc=0mm,
  boxrule=0.5pt, left=1mm, right=1mm, top=0mm, bottom=0mm,
  after skip=2mm, middle=0.3mm
}

\newcommand{\Ra}{\Rightarrow} 

\begin{document}
\begin{multicols*}{3} 

% ================= БИЛЕТ 1 =================
\begin{ticket}{1. Линейное пространство. Свойства. Примеры}
\textbf{Определение:} Множество $V$ называется линейным пространством над полем $F$ ($\mathbb{R}$ или $\mathbb{C}$), если в нем введены операции сложения и умножения на число.
Аксиомы:
\begin{enumerate}
    \item Коммутативность: $a + b = b + a$.
    \item Ассоциативность: $(a + b) + c = a + (b + c)$.
    \item Существование нуля: $\exists 0 \in V: a + 0 = a$.
    \item Существование противоположного: $\exists (-a): a + (-a) = 0$.
    \item Унитарность: $1 \cdot a = a$.
    \item Дистрибутивность: $\lambda(a + b) = \lambda a + \lambda b$ и $(\lambda + \mu)a = \lambda a + \mu a$.
    \item Ассоциативность умножения: $\lambda(\mu a) = (\lambda\mu)a$.
\end{enumerate}
\tcblower
\textbf{Свойства:}
Нулевой вектор $0$ и противоположный $(-a)$ единственны.
Умножение на ноль: $0 \cdot a = 0$. $\lambda \cdot 0 = 0$.
Умножение на $-1$: $(-1) \cdot a = -a$.
\textbf{Идеи доказательств:} $0_1 = 0_1 + 0_2 = 0_2$. $0 \cdot a = (0+0)a \Ra 0 \cdot a = 0$. $\lambda 0 = \lambda(a-a)=0$. $(-1)a+a=(-1+1)a=0 \Ra -a$.
\textbf{Примеры:} Пространство $\mathbb{R}^n$, матрицы $M_{m \times n}$.
\end{ticket}

% ================= БИЛЕТ 2 =================
\begin{ticket}{2. ЛЗ и ЛНЗ системы. Свойства. Добавление вектора}
\textbf{Определение:} Система ЛЗ, если существуют $\lambda_i$, не все равные нулю, такие что $\sum \lambda_i a_i = 0$. Иначе система линейно независима (ЛНЗ).
\textbf{Критерий:} Система ЛЗ $\Leftrightarrow$ хотя бы один вектор линейно выражается через остальные.
\textit{Идея док-ва:} $\Rightarrow$: Если $\lambda_k \neq 0$, делим на $\lambda_k$ и выражаем $a_k$. $\Leftarrow$: Переносим $a_k$ в левую часть с коэффициентом $-1 \neq 0$.
\tcblower
\textbf{Свойства:}
\begin{itemize}
    \item Содержит нулевой вектор $\Ra$ ЛЗ.
    \item Содержит ЛЗ подсистему $\Ra$ ЛЗ.
    \item Система ЛНЗ $\Ra$ любая подсистема ЛНЗ.
\end{itemize}
\textbf{О добавлении вектора:} Если система $a_1, \dots, a_k$ ЛНЗ, а с добавлением $b$ становится ЛЗ, то $b$ выражается через $a_1, \dots, a_k$.
\textit{Идея:} В нетривиальной комбинации $\sum \lambda_i a_i + \mu b = 0$ коэффициент $\mu \neq 0$, иначе исходная система была бы ЛЗ.
\end{ticket}

% ================= БИЛЕТ 3 =================
\begin{ticket}{3. Базис и размерность. Основная лемма}
\textbf{Размерность ($\dim V = n$):} В пространстве существует $n$ ЛНЗ векторов, а любая система из $n+1$ вектора ЛЗ.
\textbf{Базис:} Упорядоченная ЛНЗ система векторов, через которую линейно выражается любой вектор пространства.
\textbf{Основная лемма:} Если ЛНЗ система из $k$ векторов линейно выражается через систему из $m$ векторов, то $k \le m$.
\tcblower
\textbf{Теорема о связи:}
\begin{itemize}
    \item Если $\dim V = n$, любая система из $n$ ЛНЗ векторов --- базис.
    \item Если есть базис из $n$ векторов, то $\dim V = n$.
\end{itemize}
\textit{Идея док-ва:} Если взять $n$ ЛНЗ векторов и любой $x$, получим $n+1$ векторов, которые обязаны быть ЛЗ. По теореме о добавлении, $x$ выражается через них, значит они образуют базис.
\end{ticket}

% ================= БИЛЕТ 4 =================
\begin{ticket}{4. Координаты. Изоморфизм. Свойства}
\textbf{Координаты:} Числа $\xi_1, \dots, \xi_n$ в разложении $x = \sum \xi_i e_i$ определяются единственным образом.
\textit{Идея док-ва:} Вычитаем одно разложение из другого, в силу ЛНЗ базиса все разности коэффициентов равны нулю ($\xi_i = \eta_i$).
\textbf{Линейность:} Координаты суммы равны суммам координат, координаты произведения на число --- произведению на число.
\tcblower
\textbf{Изоморфизм ($V \cong W$):} Биекция $\varphi: V \to W$, сохраняющая аддитивность ($\varphi(x+y)=\varphi(x)+\varphi(y)$) и однородность ($\varphi(\lambda x)=\lambda\varphi(x)$).
\textbf{Свойства изоморфизма:}
\begin{itemize}
    \item $\varphi(0) = 0$ и $\varphi(-x) = -\varphi(x)$.
    \item ЛНЗ система переходит в ЛНЗ систему. \textit{Идея:} $\sum \lambda_i \varphi(x_i) = 0 \Ra \varphi(\sum \lambda_i x_i) = 0 \Ra \sum \lambda_i x_i = 0 \Ra \lambda_i = 0$.
\end{itemize}
\end{ticket}

% ================= БИЛЕТ 5 =================
\begin{ticket}{5. Теорема об изоморфизме. Координатное ЛП}
\textbf{Теорема:} Конечномерные $V$ и $W$ изоморфны тогда и только тогда, когда $\dim V = \dim W$.
\tcblower
\textbf{Идеи доказательства:}
\begin{itemize}
    \item $\Rightarrow$: Изоморфизм $\varphi$ переводит базис $V$ в базис $W$, порождая его и сохраняя линейную независимость, значит размерности совпадают.
    \item $\Leftarrow$: Если размерности равны, фиксируем базисы $E$ в $V$ и $E'$ в $W$. Строим $\varphi$, сопоставляя вектору с координатами в $E$ вектор с теми же координатами в $E'$. Оно линейно, инъективно и сюръективно.
\end{itemize}
\textbf{Следствие:} Любое $n$-мерное ЛП изоморфно координатному пространству $F^n$. Изоморфизм задается сопоставлением вектору его столбца координат.
\end{ticket}

% ================= БИЛЕТ 6 =================
\begin{ticket}{6. Подпространства. Размерность. СЛАУ}
\textbf{Определение:} Множество $L \subset V$ --- подпространство, если оно непустое ($0 \in L$) и замкнуто относительно сложения и умножения на число ($\lambda x + \mu y \in L$).
\textbf{Размерность:} $\dim L \le \dim V$. Если $\dim L = \dim V$, то $L = V$.
\textit{Идея:} Базис подпространства $L$ остается ЛНЗ в $V$, значит его длина $\le n$. Если длина $n$, то это базис всего $V$.
\tcblower
\textbf{Построение СЛАУ по базису $L$:}
Можно построить однородную СЛАУ, множеством решений которой является $L$, а ФСР --- базис $L$.
\textit{Алгоритм:} Составляем матрицу, где столбцы --- базис $L$ и произвольный вектор $x$. Приводим к ступенчатому виду. Чтобы $x \in L$, ранг не должен меняться $\Ra$ нижние элементы последнего столбца должны занулиться. Это и дает искомую СЛАУ.
\end{ticket}

% ================= БИЛЕТ 7 =================
\begin{ticket}{7. Линейные оболочки. Неполный базис}
\textbf{Линейная оболочка:} $L(a_1, \dots, a_k)$ --- множество всех линейных комбинаций этих векторов. Это подпространство.
\textbf{Размерность:} Базисом $L$ является любая максимальная ЛНЗ подсистема $\Ra \dim L = \operatorname{rang}\{a_1, \dots, a_k\}$.
\textit{Идея:} Любой вектор системы выражается через макс. ЛНЗ, значит макс. ЛНЗ порождает всю оболочку.
\tcblower
\textbf{Теорема о неполном базисе:}
ЛНЗ систему из $k$ векторов ($k < n$) можно дополнить до базиса всего пространства $V$.
\textit{Идея док-ва:} Так как $k < n$, оболочка системы не совпадает с $V$. Берем вектор $e_{k+1}$ вне оболочки и добавляем. По теореме о добавлении система остается ЛНЗ. Повторяем шаги, пока число векторов не достигнет $n$.
\end{ticket}

% ================= БИЛЕТ 8 =================
\begin{ticket}{8. Сумма и пересечение подпространств}
\textbf{Определения:} Сумма $U_1 + U_2$ --- множество векторов вида $x_1 + x_2$. Пересечение $U_1 \cap U_2$ --- векторы, принадлежащие обоим подпространствам.
\textbf{Формула Грассмана:} $\dim(U_1 + U_2) = \dim U_1 + \dim U_2 - \dim(U_1 \cap U_2)$.
\tcblower
\textit{Идея док-ва:} Берем базис $U_1 \cap U_2$, дополняем его до базиса $U_1$ и до базиса $U_2$. Объединение этих базисов порождает сумму. Линейная независимость доказывается тем, что из нулевой комбинации перенос векторов $U_2$ в одну сторону дает вектор из $U_1 \cap U_2$, откуда все коэффициенты равны нулю.
\textbf{Построение:} Базис суммы --- найти ЛНЗ систему из объединения базисов (через ступенчатый вид). Пересечение --- приравнять разложения $x$ и решить СЛАУ, ФСР дает базис.
\end{ticket}

% ================= БИЛЕТ 9 =================
\begin{ticket}{9. Прямая сумма подпространств}
\textbf{Определение:} Сумма называется прямой ($U_1 \oplus U_2$), если любой вектор из нее единственным образом представляется в виде суммы $x_1 + x_2$.
\textbf{Критерий:} Сумма прямая $\Leftrightarrow U_1 \cap U_2 = \{0\}$.
\textit{Идея:} Если есть ненулевой $v$ в пересечении, $0 = v + (-v)$ дает два разложения. Если пересечение нулевое, $x_1+x_2=x'_1+x'_2 \Ra x_1-x'_1=x'_2-x_2 \in U_1 \cap U_2 \Ra =0$.
\tcblower
\textbf{Теорема о разложении:}
Если $V = U_1 \oplus \dots \oplus U_k$, то объединение базисов $U_i$ образует базис $V$.
\textit{Идея док-ва:} Порождают, так как любой $x = \sum x_i$. Линейная независимость: нулевая комбинация означает $\sum v_i = 0$. В силу единственности разложения прямой суммы, все $v_i=0$, а из ЛНЗ базисов --- все коэффициенты нулевые.
\end{ticket}

% ================= БИЛЕТ 10 =================
\begin{ticket}{10. Элементарные преобразования матриц}
\textbf{Определение:} Элементарные преобразования:
\begin{enumerate}
    \item Перестановка строк (столбцов).
    \item Умножение строки на $\lambda \neq 0$.
    \item Прибавление к одной строке другой, умноженной на число.
\end{enumerate}
\textbf{Теорема:} Каждое преобразование строк эквивалентно умножению исходной матрицы слева на матрицу элементарного преобразования $B$.
\tcblower
\textbf{Идея доказательства:} Матрицы $B$ строятся из единичной матрицы $E$:
\begin{itemize}
    \item Для перестановки $i, j$ --- меняем $i$ и $j$ строки в $E$.
    \item Для умножения $i$-й строки --- заменяем $e_{ii}$ на $\lambda$.
    \item Для прибавления --- добавляем $\lambda$ на позицию $(i, j)$.
\end{itemize}
Все матрицы $B$ квадратные и невырождены ($\det B \neq 0$).
\end{ticket}

% ================= БИЛЕТ 11 =================
\begin{ticket}{11. Теорема о размерности суммы (Формула Грассмана)}
\textbf{Формула Грассмана:} Для любых конечномерных подпространств $U_1$ и $U_2$ пространства $V$:
\[\dim(U_1 + U_2) = \dim U_1 + \dim U_2 - \dim(U_1 \cap U_2)\]
\tcblower
\textbf{Идея доказательства:} 
Пусть $\dim(U_1 \cap U_2) = m$. Выберем базис пересечения $b_1, \dots, b_m$. 
По теореме о неполном базисе дополним его до базиса $U_1$ (добавив $a_1, \dots, a_k$) и до базиса $U_2$ (добавив $c_1, \dots, c_l$).
Объединенная система $E = \{a_i, b_j, c_p\}$ порождает $U_1 + U_2$. 
\textbf{Линейная независимость $E$:} Если линейная комбинация равна нулю, переносим векторы $c_p$ в одну сторону, а $a_i, b_j$ в другую. Полученный вектор принадлежит и $U_1$, и $U_2$, то есть $U_1 \cap U_2$. Значит, он выражается только через $b_j$. Отсюда все коэффициенты при $c_p$ равны 0. Аналогично зануляются и остальные. 
Размерность суммы: $k+m+l = (k+m) + (l+m) - m$.
\end{ticket}

% ================= БИЛЕТ 12 =================
\begin{ticket}{12. Матрица перехода. Преобразование координат}
\textbf{Матрица перехода $T$:} Матрица, столбцы которой --- координаты векторов нового базиса $E'$ в старом базисе $E$. 
Блочная запись: $E' = E \cdot T$.
\textbf{Невырожденность:} Так как $E'$ --- базис, его векторы линейно независимы. Следовательно, столбцы $T$ линейно независимы, ранг $T$ максимален, и $\det T \neq 0$.
\tcblower
\textbf{Преобразование координат:} 
Если вектор $x$ имеет столбец координат $X$ в базисе $E$ и $X'$ в базисе $E'$, то:
\[X = T X' \quad \text{или} \quad X' = T^{-1} X\]
\textbf{Идея доказательства:}
Разложим вектор по двум базисам: $x = EX$ и $x = E'X'$.
Подставим $E' = ET$: $x = (ET)X' = E(TX')$. 
В силу единственности разложения по базису $E$, координатные столбцы совпадают: $X = TX'$.
\end{ticket}

% ================= БИЛЕТ 13 =================
\begin{ticket}{13. Линейный функционал. Биортогональный базис}
\textbf{Линейный функционал:} Отображение $f: V \to F$, линейное по аргументу ($f(x+y)=f(x)+f(y)$, $f(\lambda x)=\lambda f(x)$). 
Пространство всех функционалов --- сопряженное $V^*$.
\textbf{Теорема о представлении:} $f(x) = \sum_{i=1}^n d_i \xi_i$, где $d_i = f(e_i)$, а $\xi_i$ --- координаты $x$.
\tcblower
\textbf{Биортогональный базис:} Базис $G = \{g_1, \dots, g_n\}$ в $V^*$ называется биортогональным к базису $E$ в $V$, если $g_i(e_j) = \delta_{ij}$ (символ Кронекера).
\textbf{Идея доказательства (существование):}
Строим $g_i(x) = \xi_i$. Линейная независимость: из $\sum \lambda_i g_i = 0$ при подстановке $e_j$ получаем $\lambda_j = 0$. 
\textbf{Смена базиса:} Если $E' = ET$, то координаты функционала преобразуются как $F' = F \cdot T$.
\end{ticket}

% ================= БИЛЕТ 14 =================
\begin{ticket}{14. Линейный оператор. Задание и матрица}
\textbf{Определение:} Отображение $\varphi: V \to W$ называется линейным оператором, если $\varphi(x+y)=\varphi(x)+\varphi(y)$ и $\varphi(\lambda x)=\lambda\varphi(x)$.
\textbf{Свойства:} $\varphi(0) = 0$, $\varphi(-x) = -\varphi(x)$.
\textbf{Теорема о задании:} Линейный оператор однозначно задается своими значениями на базисных векторах. 
\textit{Идея:} Для $x = \sum \xi_i e_i$ образ обязан быть $\varphi(x) = \sum \xi_i \varphi(e_i)$.
\tcblower
\textbf{Матрица оператора ($A_\varphi$):} 
В паре базисов $(E, E')$ матрица $A_\varphi$ состоит из столбцов, являющихся координатами образов базисных векторов $\varphi(e_j)$ в базисе $E'$.
Обозначение: $\varphi(E) = E' \cdot A_\varphi$.
Если $V=W$, матрица квадратная, и оператор называется преобразованием пространства $V$.
\end{ticket}

% ================= БИЛЕТ 15 =================
\begin{ticket}{15. Изоморфизм $L(V,V)$ и $M_n(F)$}
\textbf{Теорема:} Отображение $\varphi \mapsto A_\varphi$ (сопоставление оператору его матрицы в фиксированном базисе) является изоморфизмом алгебр $L(V, V)$ и $M_n(F)$.
\textbf{Следствие:} Размерность пространства операторов $\dim L(V, V) = n^2$.
\tcblower
\textbf{Идея доказательства:} 
\begin{itemize}
    \item \textbf{Биекция:} По теореме о задании на базисе, для любой матрицы $A$ существует единственный оператор.
    \item \textbf{Сложение/умножение на число:} Покоординатные операции с векторами $\varphi(e_j)$ эквивалентны операциям с матрицами.
    \item \textbf{Композиция:} $(\varphi \circ \psi)(e_j) = \varphi(\sum b_{kj}e_k) = \sum b_{kj}\varphi(e_k) = \sum_i (\sum_k a_{ik}b_{kj}) e_i$. Коэффициент --- элемент произведения матриц $A_\varphi A_\psi$.
\end{itemize}
\end{ticket}

% ================= БИЛЕТ 16 =================
\begin{ticket}{16. Преобразование координат при $\varphi$. Композиция}
\textbf{О координатах образа:} Если $x$ имеет столбец координат $X$, то его образ $y = \varphi(x)$ имеет столбец координат $Y$, вычисляемый как:
\[Y = A_\varphi X\]
\textit{Идея:} $\varphi(x) = \varphi(\sum \xi_j e_j) = \sum \xi_j (\sum a_{ij} e_i) = \sum_i (\sum_j a_{ij}\xi_j) e_i$. Выражение в скобках --- матричное умножение $AX$.
\tcblower
\textbf{Единственность матрицы:} Если в одном базисе матрицы операторов совпадают, то совпадают и сами операторы.
\textbf{Матрица произведения:} Матрица композиции операторов равна произведению их матриц:
\[A_{\varphi \circ \psi} = A_\varphi \cdot A_\psi\]
(Сначала применяется $\psi$, затем $\varphi$).
\end{ticket}

% ================= БИЛЕТ 17 =================
\begin{ticket}{17. Обратный оператор. Матрица обратного}
\textbf{Определение:} Оператор $\psi$ обратен $\varphi$ ($\psi = \varphi^{-1}$), если $\varphi \circ \psi = \psi \circ \varphi = I$ (тождественный оператор).
\textbf{Свойства:}
\begin{itemize}
    \item Единственен; линеен; $(\varphi^{-1})^{-1} = \varphi$.
    \item $(\varphi \circ \psi)^{-1} = \psi^{-1} \circ \varphi^{-1}$.
\end{itemize}
\tcblower
\textbf{Матрица обратного оператора:}
Если обратимый оператор $\varphi$ имеет матрицу $A$, то матрица обратного оператора равна обратной матрице $A^{-1}$.
\textbf{Идея доказательства:} 
Из определения $\varphi \circ \varphi^{-1} = I$. 
Переходим к матрицам (используя теорему о матрице композиции): $A \cdot A_{\varphi^{-1}} = E_n$. 
По определению обратной матрицы получаем $A_{\varphi^{-1}} = A^{-1}$.
\end{ticket}

% ================= БИЛЕТ 18 =================
\begin{ticket}{18. Критерий обратимости оператора}
\textbf{Теорема:} Для $\varphi \in L(V, V)$ (где $\dim V = n$) следующие условия эквивалентны:
\begin{enumerate}
    \item $\varphi$ обратим.
    \item $\varphi$ биективен (инъекция + сюръекция).
    \item $\det A_\varphi \neq 0$ (матрица невырождена).
    \item $\ker \varphi = \{0\}$.
    \item $\operatorname{Im} \varphi = V$.
\end{enumerate}
\tcblower
\textbf{Идея доказательства (цепочка):} 
Обратим $\Ra$ Биекция (по определению композиции, дающей $I$). 
Биекция $\Ra \ker=\{0\}$ (инъективность). 
$\ker=\{0\} \Ra \dim \operatorname{Im} = n$ (по теореме о связи размерностей) $\Ra \operatorname{Im}=V$. 
$\operatorname{Im}=V \Ra$ столбцы матрицы $A_\varphi$ порождают $V$ $\Ra$ они ЛНЗ $\Ra \det A \neq 0$. 
$\det A \neq 0 \Ra \exists A^{-1} \Ra$ задает обратный оператор.
\end{ticket}

% ================= БИЛЕТ 19 =================
\begin{ticket}{19. Ядро и образ. Связь размерностей}
\textbf{Ядро ($\ker \varphi$):} Множество векторов $x$, для которых $\varphi(x) = 0$.
\textbf{Образ ($\operatorname{Im} \varphi$):} Множество векторов $y$, таких что $\exists x : y = \varphi(x)$. 
Оба множества --- подпространства в $V$.
\textbf{Теорема о связи размерностей:} $\dim \ker \varphi + \dim \operatorname{Im} \varphi = n$.
\tcblower
\textbf{Идея доказательства:} 
Образ порождается векторами $\varphi(e_1), \dots, \varphi(e_n)$, то есть столбцами $A_\varphi$. Поэтому $\dim \operatorname{Im} \varphi = \operatorname{rang} A_\varphi$.
Ядро --- это пространство решений однородной системы $A_\varphi X = 0$. Его размерность равна $n - \operatorname{rang} A_\varphi$. 
Сумма размерностей равна $n$.
\textbf{Следствие:} В конечномерном ЛП инъективность ($\ker=\{0\}$) эквивалентна сюръективности ($\operatorname{Im}=V$).
\end{ticket}

% ================= БИЛЕТ 20 =================
\begin{ticket}{20. Смена базиса. Характеристический многочлен}
\textbf{Теорема о преобразовании матрицы:} При переходе к новому базису $E' = E \cdot T$ матрица оператора меняется по правилу:
\[A_{\varphi'} = T^{-1} A_\varphi T\]
\textit{Идея:} Из $\varphi(E') = \varphi(ET) = EA_\varphi T$ и $\varphi(E') = E'A_{\varphi'} = ETA_{\varphi'}$ следует $A_\varphi T = T A_{\varphi'}$.
\tcblower
\textbf{Характеристический многочлен:} $P_\varphi(\lambda) = \det(A_\varphi - \lambda E_n)$.
\textbf{Инвариантность:} $P_\varphi(\lambda)$ не зависит от выбора базиса.
\textbf{Идея доказательства:} 
$\det(A_{\varphi'} - \lambda E_n) = \det(T^{-1}A_\varphi T - T^{-1}(\lambda E_n)T) = \det(T^{-1}(A_\varphi - \lambda E_n)T)$.
По свойству определителя произведения: $\det(T^{-1}) \cdot \det(A_\varphi - \lambda E_n) \cdot \det(T) = \det(A_\varphi - \lambda E_n)$.
\end{ticket}

% ================= БИЛЕТ 21 =================
\begin{ticket}{21. Инвариантность ранга. Нахождение ранга}
\textbf{Теорема:} При элементарных преобразованиях строк или столбцов ранг матрицы не изменяется.
\textbf{Идея доказательства:} Достаточно показать для строк. Перестановка не меняет набор миноров. Умножение строки на $\lambda \neq 0$ умножает содержащие её миноры на $\lambda$, так что ненулевые остаются ненулевыми. Прибавление одной строки к другой сохраняет значения миноров. Максимальный порядок ненулевого минора (ранг) сохраняется.
\tcblower
\textbf{Нахождение ранга:} Любую матрицу можно привести к ступенчатому (трапециевидному) виду с помощью метода Гаусса (элементарными преобразованиями). 
В ступенчатой матрице все нулевые строки внизу, а ведущие элементы идут «лесенкой». 
\textbf{Результат:} Ранг исходной матрицы в точности равен количеству ненулевых строк в полученной ступенчатой матрице.
\end{ticket}

% ================= БИЛЕТ 22 =================
\begin{ticket}{22. Собственные числа и векторы. Спектр}
\textbf{Определения:} Ненулевой вектор $x$ --- \textbf{собственный вектор} оператора $\varphi$, если $\varphi(x) = \lambda x$. Число $\lambda$ --- \textbf{собственное число} (СЧ).
Множество всех СЧ оператора называется его \textbf{спектром}.
\textbf{Нахождение:} $\lambda$ является СЧ $\Leftrightarrow \det(A_\varphi - \lambda E) = 0$ (хар. уравнение). СВ ищутся как ФСР однородной системы $(A_\varphi - \lambda E)X = 0$.
\tcblower
\textbf{Теорема (Основное свойство):} СВ, соответствующие попарно различным СЧ, линейно независимы.
\textbf{Идея доказательства:} Индукция по числу векторов. Для $k+1$ векторов берем нулевую комбинацию $\sum c_i x_i = 0$. Применяем к ней $\varphi$, получаем $\sum c_i \lambda_i x_i = 0$. Вычитаем исходную, умноженную на $\lambda_{k+1}$, исключая $x_{k+1}$. Остается комбинация $k$ ЛНЗ векторов, откуда коэффициенты $c_i(\lambda_i - \lambda_{k+1}) = 0 \Ra c_i=0$.
\end{ticket}

% ================= БИЛЕТ 23 =================
\begin{ticket}{23. Свойства собственных векторов. Диагональность}
\textbf{Собственное подпространство $L_\lambda$:} Множество всех СВ для данного $\lambda$ и нулевого вектора. $L_\lambda = \ker(\varphi - \lambda I)$. 
\textbf{Свойство линейности:} Любая ненулевая линейная комбинация СВ для одного и того же $\lambda_0$ также является СВ с числом $\lambda_0$.
\tcblower
\textbf{Критерий диагонализируемости:} Оператор диагонализируем (его матрица имеет диагональный вид) тогда и только тогда, когда существует базис, состоящий целиком из собственных векторов. В этом базисе матрица $A_\varphi = \operatorname{diag}(\lambda_1, \dots, \lambda_n)$.
\textbf{Идея доказательства:} Если матрица диагональна, $i$-й столбец имеет лишь $\lambda_i$ на диагонали, что означает $\varphi(e_i) = \lambda_i e_i \Ra$ базис из СВ. И наоборот: если базис из СВ, координаты их образов дают диагональную матрицу.
\end{ticket}

% ================= БИЛЕТ 24 =================
\begin{ticket}{24. Кратности собственного числа}
\textbf{Алгебраическая кратность ($m_a$):} Кратность числа $\lambda$ как корня характеристического многочлена $P_\varphi(\lambda) = \det(A_\varphi - \lambda E)$.
\textbf{Геометрическая кратность ($m_g$):} Размерность собственного подпространства $L_\lambda$ ($m_g = \dim L_\lambda$).
\tcblower
\textbf{Теорема:} Для любого собственного числа $m_g \le m_a$.
\textbf{Идея доказательства:} Выберем базис из $k=m_g$ векторов в $L_\lambda$ и дополним до базиса всего пространства. В этом базисе матрица $A_\varphi$ имеет блочно-треугольный вид с верхним левым блоком $\lambda E_k$. 
Вычисляя определитель блочной матрицы $A_\varphi - \lambda E$, получаем множитель $(\lambda_0 - \lambda)^k$. Значит, корень $\lambda_0$ имеет алгебраическую кратность не менее $k$, то есть $m_a \ge m_g$.
\end{ticket}

% ================= БИЛЕТ 25 =================
\begin{ticket}{25. Прямая сумма подпространств. Базис из СВ}
\textbf{Теорема:} Сумма собственных подпространств для различных СЧ является прямой суммой: $L_{\lambda_1} \oplus \dots \oplus L_{\lambda_k}$.
\textbf{Идея:} Из равенства $x_1 + \dots + x_k = 0$ ($x_i \in L_{\lambda_i}$) в силу линейной независимости СВ для разных $\lambda_i$ следует $x_i = 0$. Пересечения нулевые $\Ra$ сумма прямая.
\tcblower
\textbf{Критерий базиса из СВ:} Базис существует $\Leftrightarrow$ 
1) Хар. многочлен распадается на линейные множители ($\sum m_a = n$). 
2) Для каждого СЧ $m_g(\lambda_i) = m_a(\lambda_i)$.
\textbf{Идея доказательства:} Размерность прямой суммы $W = \sum L_{\lambda_i}$ равна $\sum m_g$. Если $m_g = m_a$, то $\sum m_g = \sum m_a = n$. Раз $\dim W = n$, то $W$ совпадает со всем пространством, и объединение базисов подпространств $L_{\lambda_i}$ дает искомый базис из $n$ векторов.
\end{ticket}

% ================= БИЛЕТ 26 =================
\begin{ticket}{26. Билинейные формы. Разложение БФ}
\textbf{Определение:} Функция $f: V \times V \to F$, линейная по каждому аргументу. 
Виды: симметричная ($f(x,y)=f(y,x)$) и кососимметричная ($f(x,y)=-f(y,x)$).
\textbf{Полуторалинейная форма:} Над $\mathbb{C}$, линейна по первому и антилинейна по второму аргументу.
\tcblower
\textbf{Теорема о разложении:} Любая БФ единственным образом представляется суммой симм. ($f_1$) и кососимм. ($f_2$) форм:
$f_1(x,y) = \frac{1}{2}(f(x,y)+f(y,x))$, $f_2(x,y) = \frac{1}{2}(f(x,y)-f(y,x))$.
\textbf{Замечание:} Симметричная БФ однозначно определяется значениями на диагонали $f(x,x)$ через тождество поляризации: 
$f(x,y) = \frac{1}{2}(f(x+y,x+y) - f(x,x) - f(y,y))$.
\end{ticket}

% ================= БИЛЕТ 27 =================
\begin{ticket}{27. Матрица БФ. Преобразование при смене базиса}
\textbf{Матрица БФ ($A_f$):} В базисе $E$ элементы $a_{ij} = f(e_i, e_j)$. 
Значение формы: $f(x, y) = X^T A_f Y$.
\textbf{Свойства:} Отображение $f \mapsto A_f$ --- изоморфизм пространств БФ и матриц.
\tcblower
\textbf{Преобразование матрицы БФ:} При переходе к новому базису $E' = E \cdot T$ матрица БФ преобразуется по формуле:
\[A_{f'} = T^T A_f T\]
\textbf{Идея доказательства:} Столбцы координат преобразуются как $X = T X'$ и $Y = T Y'$. 
Подставим в формулу значения формы: 
$f(x, y) = (T X')^T A_f (T Y') = (X')^T (T^T A_f T) Y'$.
Так как $f(x, y) = (X')^T A_{f'} Y'$ и равенство верно для любых $X', Y'$, получаем $A_{f'} = T^T A_f T$.
\end{ticket}

% ================= БИЛЕТ 28 =================
\begin{ticket}{28. Квадратичные формы. Методы Лагранжа и Якоби}
\textbf{Квадратичная форма:} $Q(x) = f(x,x)$, где $f$ --- симметричная полярная БФ. Канонический вид --- сумма квадратов $\sum \alpha_i \xi_i^2$.
\textbf{Метод Лагранжа:} Алгоритм приведения к каноническому виду выделением полных квадратов.
\textit{Идея:} Если $\exists a_{11} \neq 0$, все слагаемые с $\xi_1$ собираем в полный квадрат, делаем замену и переходим к $n-1$ переменным. Если все $a_{ii}=0$, но $a_{12} \neq 0$, делаем невырожденную замену $\xi_1 = \eta_1+\eta_2, \xi_2 = \eta_1-\eta_2$, получая квадраты $\eta_1^2, \eta_2^2$, и сводим к первому случаю.
\tcblower
\textbf{Метод Якоби:} Если все главные угловые миноры матрицы формы $\Delta_k \neq 0$, существует преобразование к виду:
$Q(x) = \frac{\Delta_1}{\Delta_0} \eta_1^2 + \frac{\Delta_2}{\Delta_1} \eta_2^2 + \dots + \frac{\Delta_n}{\Delta_{n-1}} \eta_n^2$, где $\Delta_0 = 1$.
\end{ticket}

% ================= БИЛЕТ 29 =================
\begin{ticket}{29. Закон инерции квадратичных форм}
\textbf{Нормальный вид:} Над $\mathbb{R}$ форма приводится к сумме квадратов с коэффициентами $\pm 1$:
$Q(x) = \xi_1^2 + \dots + \xi_p^2 - \xi_{p+1}^2 - \dots - \xi_{p+q}^2$.
Здесь $p$ --- положительный, а $q$ --- отрицательный индекс инерции.
\tcblower
\textbf{Теорема (Закон инерции):} Индексы $p$ и $q$ (и дефект) инвариантны и не зависят от выбора нормального базиса.
\textbf{Идея доказательства:} От противного. Пусть в базисе $E$ инд. $p$, а в $E'$ инд. $r$, и $p > r$. Строим $L_1$ на первых $p$ векторах $E$ и $L_2$ на последних $n-r$ векторах $E'$. 
$\dim(L_1 \cap L_2) \ge p + (n-r) - n = p - r > 0$. 
Берем ненулевой вектор $x_0 \in L_1 \cap L_2$. С одной стороны ($L_1$), форма на нем $>0$. С другой ($L_2$) --- $\le 0$. Получили противоречие, значит $p = r$.
\end{ticket}

% ================= БИЛЕТ 30 =================
\begin{ticket}{30. Классификация КФ. Критерий Сильвестра}
\textbf{Классификация:} В зависимости от знака $f(x,x)$ формы бывают: 
положительно/отрицательно (полу)определенные и незнакоопределенные.
\tcblower
\textbf{Критерий Сильвестра:} 
\begin{enumerate}
    \item Форма положительно определена $\Leftrightarrow$ все $\Delta_k > 0$.
    \item Форма отрицательно определена $\Leftrightarrow$ знаки $\Delta_k$ чередуются ($\Delta_1 < 0, \Delta_2 > 0, \dots$).
\end{enumerate}
\textbf{Идея доказательства:} $\Rightarrow$: Если $\Delta_k = 0$, однородная система из первых $k$ уравнений имеет решение $x_0 \neq 0$, на котором форма $f(x_0, x_0) = 0$, что противоречит строгой $>0$. Раз $\Delta_k \neq 0$, по методу Якоби все коэффициенты $\frac{\Delta_k}{\Delta_{k-1}} > 0$, откуда $\Delta_k > 0$.
$\Leftarrow$: Если $\Delta_k > 0$, по методу Якоби все коэффициенты положительны, форма есть сумма квадратов $\Ra f > 0$. Для отриц. определенной применяем рассуждения к форме $-f$.
\end{ticket}

% ================= БИЛЕТ 31 =================
\begin{ticket}{31. Нахождение обратной матрицы методом эл. преобразований}
\textbf{Теорема:} Квадратная матрица $A$ путем элементарных преобразований строк может быть приведена к единичной матрице $E$ $\Leftrightarrow$ $\det A \neq 0$.
\textbf{Идея доказательства:} 
$\Rightarrow$: $B_k \dots B_1 A = E \Ra \det(B_k)\dots\det(A) = 1 \Ra \det A \neq 0$. 
$\Leftarrow$: Если $\det A \neq 0$, то $\operatorname{rang} A = n$. Ступенчатая форма не имеет нулевых строк. Обратным ходом обнуляем элементы над диагональю и получаем $E$.
\tcblower
\textbf{Теорема (Алгоритм):} Если к расширенной матрице $(A \mid E)$ применить элементарные преобразования строк, приводящие $A$ к $E$, то она преобразуется к виду $(E \mid A^{-1})$.
\textbf{Идея доказательства:} 
Приведение $A$ к $E$ означает умножение слева на матрицу $B$ (произведение матриц эл. преобразований), такую что $BA = E$. Значит, $B = A^{-1}$. 
Умножая расширенную матрицу на $B$, получаем:
\[B \cdot (A \mid E) = (BA \mid BE) = (E \mid A^{-1})\]
\end{ticket}

% ================= БИЛЕТ 32 =================
\begin{ticket}{32. Евклидовы и унитарные пространства. Норма}
\textbf{Евклидово ЛП:} Пространство над $\mathbb{R}$ со скалярным произведением $(x, y) \in \mathbb{R}$, удовлетворяющим аксиомам:
1. Симметричность: $(x, y) = (y, x)$.
2. Аддитивность: $(x_1 + x_2, y) = (x_1, y) + (x_2, y)$.
3. Однородность: $(\lambda x, y) = \lambda (x, y)$.
4. Положительная определенность: $(x, x) > 0$ при $x \neq 0$, и $(x, x) = 0 \iff x = 0$.
\tcblower
\textbf{Унитарное ЛП:} Пространство над $\mathbb{C}$. Скалярное произведение $(x, y) \in \mathbb{C}$ удовлетворяет:
1. Эрмитова симметричность: $(x, y) = \overline{(y, x)}$.
Аксиомы 2-4 аналогичны, но $\lambda \in \mathbb{C}$, а $(x, x) \in \mathbb{R}$.
\textbf{Норма вектора (длина):} $\|x\| = \sqrt{(x, x)}$. 
Норма неотрицательна и равна нулю только для $x = 0$.
\end{ticket}

% ================= БИЛЕТ 33 =================
\begin{ticket}{33. Неравенства Коши-Буняковского и треугольника}
\textbf{Теорема Коши-Буняковского:} $|(x, y)| \le \|x\| \cdot \|y\|$.
\textbf{Идея доказательства:} Рассмотрим вектор $z = x - \lambda y$. В силу аксиом, $(z, z) \ge 0$. Подставляя $\lambda = \frac{(x, y)}{(y, y)}$ (если $y \neq 0$) в неравенство $(x - \lambda y, x - \lambda y) \ge 0$ и раскрывая скобки (учитывая сопряжение для унитарного пространства), получаем:
\[\|x\|^2 - \frac{|(x, y)|^2}{\|y\|^2} \ge 0 \Ra |(x, y)|^2 \le \|x\|^2 \|y\|^2\]
\tcblower
\textbf{Неравенство треугольника:} $\|x + y\| \le \|x\| + \|y\|$.
\textbf{Идея доказательства:} 
Возведем в квадрат: $\|x + y\|^2 = (x + y, x + y) = \|x\|^2 + (x, y) + (y, x) + \|y\|^2$.
Сумма $(x, y) + (y, x) = 2 \operatorname{Re}(x, y) \le 2|(x, y)|$. 
Применяя неравенство Коши-Буняковского:
\[\|x + y\|^2 \le \|x\|^2 + 2\|x\|\|y\| + \|y\|^2 = (\|x\| + \|y\|)^2\]
Извлекаем корень $\Ra$ требуемое неравенство.
\end{ticket}

% ================= БИЛЕТ 34 =================
\begin{ticket}{34. СЛАУ. Основные определения}
\textbf{СЛАУ:} Система $m$ линейных уравнений с $n$ неизвестными:
\[\sum_{j=1}^n a_{ij}x_j = b_i, \quad i=1,\dots,m\]
\textbf{Матрицы:} $A$ --- матрица системы (из коэффициентов $a_{ij}$), $\overline{A}$ --- расширенная матрица (добавлен столбец свободных членов $b_i$).
\tcblower
\textbf{Виды систем:}
\begin{itemize}
    \item \textbf{Однородная}, если все $b_i = 0$. \textbf{Неоднородная}, если $\exists b_i \neq 0$.
    \item \textbf{Совместная}, если есть хотя бы 1 решение. \textbf{Несовместная} --- нет решений.
    \item \textbf{Определенная}, если решение единственно. \textbf{Неопределенная} --- более одного решения (бесконечно много).
\end{itemize}
\textbf{Общее решение:} Совокупность всех частных решений системы.
\end{ticket}

% ================= БИЛЕТ 35 =================
\begin{ticket}{35. Теорема Крамера}
\textbf{Теорема:} Если в СЛАУ число уравнений равно числу неизвестных ($m=n$) и $\det A \neq 0$, то система имеет единственное решение, вычисляемое по формулам:
\[x_k = \frac{\Delta_k}{\Delta}, \quad k = 1, \dots, n\]
где $\Delta = \det A$, а $\Delta_k$ --- определитель матрицы, полученной из $A$ заменой $k$-го столбца на столбец свободных членов $b$.
\tcblower
\textbf{Идея доказательства:} 
1. \textbf{Существование:} $\det A \neq 0 \Ra \exists A^{-1}$. Умножим $Ax = b$ слева на $A^{-1} \Ra x = A^{-1}b$.
2. \textbf{Единственность:} Если есть решения $x^{(1)}$ и $x^{(2)}$, то $A(x^{(1)} - x^{(2)}) = 0 \Ra x^{(1)} = x^{(2)}$.
3. \textbf{Формулы:} Из $x = A^{-1}b$ получаем $x_k = \frac{1}{\det A} \sum_{j=1}^n A_{jk} b_j$, где $A_{jk}$ --- алгебраическое дополнение. Сумма $\sum A_{jk} b_j$ является разложением определителя $\Delta_k$ по $k$-му столбцу.
\end{ticket}

% ================= БИЛЕТ 36 =================
\begin{ticket}{36. Метод Гаусса. Теорема Кронекера-Капелли}
\textbf{Метод Гаусса:} Алгоритм приведения матрицы к ступенчатому (трапециевидному) виду элементарными преобразованиями. Прямой ход (исключение переменных) дает ступенчатую форму, обратный --- решение.
\textbf{Эквивалентность:} Элементарные преобразования (перестановка уравнений, умножение на число, прибавление другого уравнения) не меняют множество решений.
\tcblower
\textbf{Теорема Кронекера-Капелли:} СЛАУ совместна $\Leftrightarrow \operatorname{rang} A = \operatorname{rang} \overline{A}$.
\textbf{Идея доказательства:} 
$\Rightarrow$: Если решение есть, столбец $b$ есть линейная комбинация столбцов $A$, добавление $b$ не увеличивает ранг.
$\Leftarrow$: Если ранги равны, то при приведении $\overline{A}$ к ступенчатому виду не возникает строк вида $(0 \dots 0 \mid c)$ с $c \neq 0$. Значит, система разрешима обратным ходом метода Гаусса.
\textbf{Следствия:} Если $\operatorname{rang} A = n$ --- определена. Если $<n$ --- неопределена.
\end{ticket}

% ================= БИЛЕТ 37 =================
\begin{ticket}{37. Однородные системы. Свойства решений}
\textbf{Однородная система:} $Ax = 0$. Она всегда совместна, так как всегда есть \textbf{тривиальное} (нулевое) решение $x = 0$.
\textbf{Свойства решений:} Множество решений образует линейное подпространство. Если $x^{(1)}$ и $x^{(2)}$ --- решения, то:
1. $A(x^{(1)} + x^{(2)}) = Ax^{(1)} + Ax^{(2)} = 0 \Ra$ сумма является решением.
2. $A(\lambda x^{(1)}) = \lambda A x^{(1)} = 0 \Ra$ произведение на число является решением.
\tcblower
\textbf{Критерий нетривиальных решений:}
Однородная система имеет нетривиальное (ненулевое) решение $\Leftrightarrow \operatorname{rang} A < n$ (число неизвестных).
Для квадратной матрицы ($m=n$) это равносильно условию $\det A = 0$.
\textbf{Идея:} Нетривиальные решения существуют тогда и только тогда, когда в ступенчатом виде системы остается хотя бы одна свободная переменная, что бывает при ранге строго меньшем $n$.
\end{ticket}

% ================= БИЛЕТ 38 =================
\begin{ticket}{38. ФСР однородной системы и структура общего решения}
\textbf{ФСР:} Фундаментальная система решений --- любой базис в подпространстве решений однородной системы $Ax = 0$. Состоит из ЛНЗ решений $y^{(1)}, \dots, y^{(k)}$, через которые выражается любое решение.
\textbf{Теорема о ФСР:} Если $\operatorname{rang} A = r < n$, то ФСР существует и состоит ровно из $n - r$ решений.
\tcblower
\textbf{Структура общего решения:}
\[X_{оо} = c_1 y^{(1)} + c_2 y^{(2)} + \dots + c_{n-r} y^{(n-r)}\]
\textbf{Идея построения:} 
Приводим $A$ к ступенчатому виду. Выделяем $r$ главных и $n-r$ свободных переменных. Задавая свободным переменным значения строк единичной матрицы порядка $n-r$, находим $n-r$ решений. Они ЛНЗ (т.к. блок свободных переменных единичный) и любое решение выражается через них (т.к. решение однозначно определяется свободными переменными).
\end{ticket}

% ================= БИЛЕТ 39 =================
\begin{ticket}{39. Неоднородные системы. Теорема об общем решении}
	\textbf{Теорема:} Общее решение $X$ совместной неоднородной системы $Ax = b$ равно сумме любого частного решения этой системы $X_{ch}|$ и общего решения соответствующей однородной системы $X_{oo}$ ($Ax=0$):
	\[X = X_{oo} + X_{ch}\]
\tcblower
\textbf{Идея доказательства:} 
\begin{itemize}
	\item \textbf{Сумма является решением:} $A(X_{oo} + X_{ch}) = AX_{oo} + AX_{ch} = 0 + b = b$. Следовательно, вектор $X_{oo} + X_{ch}$ удовлетворяет исходной системе.
	\item \textbf{Любое решение имеет такой вид:} Пусть $X$ --- произвольное решение $Ax=b$. Рассмотрим вектор $y = X - X_{ch}$. $Ay = AX - AX_{ch} = b - b = 0$. Значит, $y$ --- решение однородной системы ($y = X_{oo}$). Отсюда $X = X_{oo} + X_{ch}$.
\end{itemize}
Следствие: Общее решение зависит от $n-r$ произвольных постоянных.
\end{ticket}

% ================= БИЛЕТ 40 =================
\begin{ticket}{40. Основная лемма о линейной зависимости}
\textbf{Лемма:} Пусть система $S_1 = \{a_1, \dots, a_k\}$ ЛНЗ, и каждый ее вектор линейно выражается через систему $S_2 = \{b_1, \dots, b_m\}$. Тогда $k \le m$.
\textbf{Идея доказательства (от противного):} 
Предположим, $k > m$. Выразим каждый вектор $a_j = \sum_{i=1}^m \alpha_{ij} b_i$. 
Рассмотрим нулевую линейную комбинацию: $\sum_{j=1}^k x_j a_j = 0$.
\tcblower
Подставим выражения для $a_j$:
\[\sum_{j=1}^k x_j \left( \sum_{i=1}^m \alpha_{ij}b_i \right) = \sum_{i=1}^m \left( \sum_{j=1}^k \alpha_{ij}x_j \right) b_i = 0\]
Приравняв коэффициенты при $b_i$ к нулю, получим однородную СЛАУ из $m$ уравнений с $k$ неизвестными ($x_1, \dots, x_k$).
Так как $k > m$, ранг системы $\le m < k$, следовательно, система имеет \textbf{нетривиальное решение}.
Это означает, что исходная линейная комбинация $a_j$ с этими нетривиальными $x_j$ дает $0$, что противоречит ЛНЗ системы $S_1$. Значит, $k \le m$.
\end{ticket}

\end{multicols*}
\end{document}
